\documentclass[../../main.tex]{subfiles}
\begin{document}
% SOURCE p=245; printed=230; object=ch07-finite-and-infinite-products-transition; region=upper
Radon 乘积的概念直接推广到任意有限个因子。更有意思的是，若各空间紧，且各测度归一化为总质量 $1$，理论还可推广到无限多个因子。
% SUPPLEMENT: id=v14-ch07-finite-products-details; parent=ch07-finite-and-infinite-products-transition; kind=omitted-proof+motivation+hypotheses; origin=independent-derivation
\begin{remark}
\textbf{有限乘积的括号与次序（重构新增）。}逐次应用两因子的定义可构造任意有限个 Radon 测度的乘积，且所得空间仍 LCH。它不依赖括号，在交换坐标的自然同胚下也不依赖因子次序。先对有限个分离变量函数 $f=\prod_{j=1}^ng_j(x_j)$、$g_j\in C_c(X_j)$ 检验：反复用紧支撑函数的累次积分公式，任意括号均给 $\int f=\prod_j\int g_j$，各积分有限。

这些函数的有限线性组合可在一个固定紧容器内一致逼近任意 $f\in C_c(\prod_jX_j)$。具体地，取各支集投影的相对紧开邻域 $U_j$，在紧乘积 $\prod_j\overline U_j$ 上用有限张量代数：它含 $1$、对共轭封闭，并由某个不同坐标的连续函数分离两点，所以 Stone--Weierstrass 给一致逼近。再乘以各坐标上在支集投影等于 $1$ 的紧支撑截断，恰如\cref{proposition:ch07-tensor-density}，得到支撑在固定紧集 $K=\prod_j\operatorname{supp}\phi_j$ 内的近似。任意两种括号得到的 Radon 测度 $\rho,\sigma$ 都在 $K$ 上有限，因而
\begin{align*}
\left|\int f\,\mathrm d\rho-\int f\,\mathrm d\sigma\right|
\leqslant(\rho(K)+\sigma(K))\|f-F\|_u
\end{align*}
对每个张量和近似 $F$ 成立。误差趋零，得到全部 $C_c$ 积分一致，再由 Riesz 唯一性得到测度相同。置换坐标时，把一方测度沿自然同胚推送；同胚保持紧集、开集和各项 Radon 条件，相同论证成立。

各因子紧且总质量为 $1$ 时，常函数 $1$ 属于 $C_c$，上面的张量积分公式说明有限乘积总质量也为 $1$。归一化使增添一个未被函数使用的坐标不会改变积分，这是下一个无限乘积构造相容性的关键。
\end{remark}
% END SUPPLEMENT: id=v14-ch07-finite-products-details
% SOURCE p=245; printed=230; object=ch07-cylinder-projection-definition; region=upper-middle
确切地说，设 $\{X_\alpha\}_{\alpha\in A}$ 为紧 Hausdorff 空间族，对每个 $\alpha$，$\mu_\alpha$ 为 $X_\alpha$ 上满足 $\mu_\alpha(X_\alpha)=1$ 的 Radon 测度。由 Tychonoff 定理，$X=\prod_{\alpha\in A}X_\alpha$ 为紧 Hausdorff 空间。我们希望在 $X$ 上定义 Radon 测度 $\mu$，使每个 $E_\alpha\in\mathscr B_{X_\alpha}$ 且除有限多个 $\alpha$ 外均有 $E_\alpha=X_\alpha$ 时，满足
\begin{align*}
\mu\left(\prod_{\alpha\in A}E_\alpha\right)=\prod_{\alpha\in A}\mu_\alpha(E_\alpha).
\end{align*}
右侧除有限项外都为 $1$，所以有良好定义。
\begin{definition}[有限坐标投影]\label{definition:ch07-finite-coordinate-projection}
给定互不相同的 $\alpha_1,\ldots,\alpha_n\in A$，记 $\pi_{(\alpha_1,\ldots,\alpha_n)}$ 为 $X$ 到 $\prod_{j=1}^nX_{\alpha_j}$ 的自然投影：
\begin{align*}
\pi_{(\alpha_1,\ldots,\alpha_n)}(x)=(x_{\alpha_1},\ldots,x_{\alpha_n}).
\end{align*}
\end{definition}
于是 $\pi_{(\alpha_1,\ldots,\alpha_n)}^{-1}(E_{\alpha_1}\times\cdots\times E_{\alpha_n})=\prod_{\alpha\in A}E_\alpha$，其中未列出的坐标取 $E_\alpha=X_\alpha$。
% SUPPLEMENT: id=v14-ch07-cylinder-projection-details; parent=ch07-cylinder-projection-definition; kind=motivation+hypotheses+source-clarification; origin=independent-derivation
\begin{remark}
\textbf{有限坐标与归一化（重构新增）。}这些柱集只规定有限个坐标，可先用有限乘积确定其概率，再要求所有有限描述相容。每个因子的质量为 $1$，所以各因子非空；沿用 Tychonoff 定理的选择公理背景，自然投影满射。这里把原文隐含的“所列坐标互不相同”明确写出。若重复一个下标，例如 $(\alpha,\alpha)$，其投影集中在对角线上，边缘测度一般不是两个独立副本的乘积；应先删除重复坐标，而不能使用下一定理的独立乘积公式。
\end{remark}
% END SUPPLEMENT: id=v14-ch07-cylinder-projection-details
% SOURCE p=245; printed=230; object=ch07-product-regular-borel; number=7.28; region=middle
\begin{theorem}\label{theorem:ch07-product-regular-borel}
设对每个 $\alpha\in A$，$\mu_\alpha$ 是紧 Hausdorff 空间 $X_\alpha$ 上满足 $\mu_\alpha(X_\alpha)=1$ 的 Radon 测度。则 $X=\prod_{\alpha\in A}X_\alpha$ 上存在唯一 Radon 测度 $\mu$，使对任意有限个互不相同的 $\alpha_1,\ldots,\alpha_n\in A$ 及 $\prod_{j=1}^nX_{\alpha_j}$ 中的任意 Borel 集 $E$，有
\begin{align*}
\mu\bigl(\pi_{(\alpha_1,\ldots,\alpha_n)}^{-1}(E)\bigr)
=(\mu_{\alpha_1}\mathbin{\widehat\times}\cdots\mathbin{\widehat\times}\mu_{\alpha_n})(E).
\end{align*}
\end{theorem}
% SOURCE p=245; printed=230; object=ch07-finite-coordinate-functions-definition; region=middle-lower
\begin{definition}[有限坐标连续函数]\label{definition:ch07-finite-coordinate-functions}
以 $C_F(X)$ 表示仅依赖有限多个坐标的连续函数所成集合，即具有形式 $f=g\circ\pi_{(\alpha_1,\ldots,\alpha_n)}$ 的 $f\in C(X)$，其中 $g\in C(\prod_{j=1}^nX_{\alpha_j})$。
\end{definition}
% SOURCE p=245,246; printed=230-231; object=ch07-product-regular-borel-proof; region=lower-upper
\begin{proof}
对这样的函数定义
\begin{align*}
I(f)=\int g\,\mathrm d(\mu_{\alpha_1}\mathbin{\widehat\times}\cdots\mathbin{\widehat\times}\mu_{\alpha_n}).
\end{align*}
给坐标集增加若干额外坐标不会改变这个公式，因为全部 $\mu_\alpha(X_\alpha)=1$。所以 $I$ 是 $C_F(X)$ 上良定义的正线性泛函，且 $|I(f)|\leqslant\|f\|_u$；常函数时取等号。

$C_F(X)$ 显然是分离点、包含常函数并对复共轭封闭的代数，由 Stone--Weierstrass 定理，它在 $C(X)$ 中稠密。因此 $I$ 唯一延拓为 $C(X)$ 上范数为 $1$ 的正线性泛函，Riesz 定理产生唯一 Radon 测度 $\mu$，使每个 $f\in C_F(X)$ 都有 $I(f)=\int f\,\mathrm d\mu$。

给定这些坐标，令 $\mu_{(\alpha_1,\ldots,\alpha_n)}=\mu\circ\pi_{(\alpha_1,\ldots,\alpha_n)}^{-1}$。它是有限乘积空间上的 Borel 测度，满足
\begin{align*}
\int g\,\mathrm d\mu_{(\alpha_1,\ldots,\alpha_n)}
=\int g\circ\pi_{(\alpha_1,\ldots,\alpha_n)}\,\mathrm d\mu
\end{align*}
先对 Borel 集示性函数成立，进而由通常的线性和逼近论证，对任意有界 Borel 函数成立。特别地，由 $\mu$ 的定义，对每个连续 $g$ 有
\begin{align*}
\int g\,\mathrm d\mu_{(\alpha_1,\ldots,\alpha_n)}
=\int g\,\mathrm d(\mu_{\alpha_1}\mathbin{\widehat\times}\cdots\mathbin{\widehat\times}\mu_{\alpha_n}).
\end{align*}
只需证明推送测度 Radon，Riesz 的唯一性便给出它与有限 Radon 乘积相同。

简记投影为 $\pi$。若 $E$ 是有限乘积中的 Borel 集，$\mu$ 的正则性使对任意 $\varepsilon>0$，存在紧集 $K\subseteq\pi^{-1}(E)$，使 $\mu(K)>\mu(\pi^{-1}(E))-\varepsilon$。$K'=\pi(K)$ 为 $E$ 的紧子集，且
\begin{align*}
\mu_{(\alpha_1,\ldots,\alpha_n)}(K')
=\mu(\pi^{-1}(K'))\geqslant\mu(K)
>\mu_{(\alpha_1,\ldots,\alpha_n)}(E)-\varepsilon.
\end{align*}
所以推送测度内正则；对 $E^c$ 使用同一论证又得到外正则。因而它 Radon，证明完成。
\end{proof}
% SUPPLEMENT: id=v14-ch07-infinite-products-details; parent=ch07-product-regular-borel; kind=proof-detail+motivation+hypotheses; origin=independent-derivation
\begin{remark}
\textbf{相容性、正延拓、推送与唯一性（重构新增）。}先核验 $I$ 的表示独立性。若 $f$ 同时依赖有限坐标集 $B,C$，把两个表示都提到 $D=B\cup C$。投影 $X\to\prod_{\alpha\in D}X_\alpha$ 满射，故两提升函数处处相等。有限乘积中，每添加一个未使用的坐标，其积分乘上该因子的总质量 $1$，所以两次积分相同。线性组合可在共同有限坐标集内计算，得线性；若 $f\geqslant0$，满射性使其代表 $g\geqslant0$，所以 $I(f)\geqslant0$。同样 $\|g\|_u=\|f\|_u$，有限乘积为概率测度，给 $|I(f)|\leqslant\|f\|_u$，而 $I(1)=1$ 给范数恰为 $1$。

有限坐标函数的和、积依赖下标并集，共轭依赖同一坐标，所以它是自共轭代数。不同的 $x,y\in X$ 至少在一个坐标 $\alpha$ 不同；紧 Hausdorff 空间 $X_\alpha$ 上有连续函数分离那两个坐标值，复合投影便在 $C_F(X)$ 中分离 $x,y$。它含常数，因此 Stone--Weierstrass 的全部条件成立。有界性给唯一连续延拓。还要核验延拓保持正性：若 $f\geqslant0$，取 $h_n\in C_F$ 一致趋于 $f$，取实部 $u_n$ 并置 $v_n=u_n+\|u_n-f\|_u\cdot1$，则 $v_n\geqslant0$、$v_n\in C_F$ 且 $v_n\to f$ 一致。因此延拓值为 $\lim I(v_n)\geqslant0$。所得 $\mu(X)=I(1)=1$，所以有限 Radon 测度的全体内正则性可用。

推送积分恒等式对示性函数直接来自定义，对非负简单函数来自线性性，对非负 Borel 函数来自逐增简单逼近及两边的单调收敛；有界复 Borel 函数再分实、虚正负四部分，因为总质量有限，各积分都有限。推送的外正则性也可写清：令 $\lambda=\mu\circ\pi^{-1}$，对 $E^c$ 取紧 $L\subseteq E^c$，使 $\lambda(E^c\setminus L)<\varepsilon$；有限乘积 Hausdorff，$L$ 闭，所以 $O=L^c$ 开且含 $E$，并有 $\lambda(O\setminus E)<\varepsilon$。推送总质量为 $1$，故紧集有限，连同内、外正则性确为 Radon。此后才能合法用 Riesz 唯一性识别有限维边缘。

若另一个 Radon 测度 $\eta$ 也具有这些边缘分布，则柱函数的积分恒等式给 $\int f\,\mathrm d\eta=I(f)$（$f\in C_F$）；特别地总质量为 $1$，所以两积分泛函连续。由稠密性在全部 $C(X)$ 上相同，Riesz 唯一性给 $\eta=\mu$，完成题目要求的唯一性。若指标集为空，乘积约定为单点空间，取其单位点质量；空坐标边缘条件也规定总质量为 $1$。

动机是先规定所有有限实验的联合分布，再由紧性和连续函数逼近恢复整个无限乘积上的测度。紧 Hausdorff 条件保证 Tychonoff 紧性、连续分离及 Riesz 表示；归一化为 $1$ 保证增加闲置坐标不改变积分；Radon 条件使有限维连续读数能唯一确定全部 Borel 质量。下标族可以任意大，证明从未要求它可数。
\end{remark}
% END SUPPLEMENT: id=v14-ch07-infinite-products-details
\subfile{section-04-exercises.tex}
\end{document}
